Training-free caching accelerates a video diffusion model by evaluating a subset of the denoising steps and reusing the last exact output elsewhere. Existing methods decide which steps to evaluate with a threshold on a quantity that has already been observed, and spend their design effort on how a skipped step is reconstructed. We measure what a skipped step actually costs, and the cost turns out to be positional. On CogVideoX-2B at thirty steps, skipping the first step costs 0.57 LPIPS against the full run while every step after the tenth costs about 0.04, so the first ten steps carry about three quarters of the sensitivity, and two prompts agree on the order at 0.94 correlation. Because the profile is a property of the sampler, \method{} can place the calls by solving the induced allocation program; the program admits a greedy optimum, and because the profile does not depend on the trajectory, the model-predictive loop collapses to an open-loop schedule with an exact head and gaps that grow along the tail, whose number of calls matches the budget by construction. Evaluated inside one released pipeline at matched calls, \method{} reaches 0.151 LPIPS where the risk-based cache of RA-CFGCache reaches 0.169 on CogVideoX-2B, 0.121 against 0.132 at 49 frames, and 0.134 against 0.154 on Wan2.1. Code, schedules and all evaluation traces are released at \url{https://github.com/Yizhou-Yang/MPCache}.
